% p3-10-worked-examples

\section{P3 §10. Worked Examples}

10. Worked Examples

           Figure (fig-p3-ch10-examples). Worked-examples triptych --- phi / Lucas L7-L8 / AlphaPhi
                                              disambiguation.

  Preamble. All examples produce mathematical approximations. No example constitutes
  evidence that a physical constant is structurally determined by the GOLDEN BRIDGE (Trinity-Pellis Bridge)
  framework. Physical constants are governed by physical laws, not compression
  grammars. The selection of physical constants as worked examples reflects only that
  they are familiar, precisely measured real numbers to which any symbolic
  approximation scheme can be applied.

  10.1 The Golden Ratio $\varphi
  x = $\varphi$ approx 1.61803398874989.

  Trinity representation: T(1, 0, 1, 0, 0) = $\varphi$. Exact; MDL score = ell((1,0,1,0,0)) = 1 + 1 +
  2 + 1 + 1 = 6 bits (1 for n=1, 1 for k=0, 2 for p=1, 1 for m=0, 1 for q=0).

  Pellis expansion: $\varphi$ $\in$ Z[$\varphi$] (with a = 0, b = 1). By Corollary 4.5, terminates at depth 1:
  c0 = 1, c1 = 1, remainder 0.

  Phase transition exponent: muT($\varphi$) = 0 (maximally Trinity-dominant; Proposition 8.2).

  Comment. The golden ratio is simultaneously the simplest Trinity monomial T(1,0,1,0,0)
  = $\varphi$ and the simplest element of Z[$\varphi$] outside the integers. It serves as the ground state
  of the framework.

  10.2 Lucas Seed L7 = 29 and L8 = 47
  Trinity representation: T(29, 0, 0, 0, 0) = 29. MDL score = lceil log2(30)rceil = 5 bits.

  Pellis expansion: Terminates exactly at depth 0: c0 = 29, remainder 0. muT(29) = 0.

        Observation. L7 = 29 and L8 = 47 satisfy the Lucas recurrence Ln+2 = Ln+1 +
        Ln. In Trinity monomial terms, neither 29 nor 47 simplifies to a product of the
        transcendental generators at low complexity; both are trivially integer monomials

        with the highest possible Trinity-dominance score (muT = 0). The interesting
        structure of Lucas numbers under Trinity representation emerges at larger indices
        where Ln approx $\varphi$n becomes a more accurate approximation, i.e., where T(1, 0, n,
        0, 0) = $\varphi$n provides a non-trivial Trinity approximant.

  10.3 Fibonacci Seed F17 = 1597
  x = F17 = 1597.

  Trinity representation: T(1597, 0, 0, 0, 0) = 1597. MDL score = lceillog2(1598)rceil =
  11 bits.

  Pellis expansion: Exact termination (Corollary 4.5): c0 = 1597, remainder 0.

  Alternative Trinity approximant: Binet's formula gives F17=($\varphi$17-psi17)/5, so the
  leading term $\varphi$17/5approx 1596.999875 differs from the exact integer by about 1.25$\times
  10-4. This confirms that for Fibonacci and Lucas ring elements, the exact integer
  representation remains MDL-optimal unless the conjugate correction term -psi17/5 is
  included.

  10.4 The Number pi
  x = pi approx 3.14159265358979.

  Trinity representation at L = 4: T(1,0,0,1,0) = pi. Exact; MDL score = 4 bits.

  Pellis expansion: Leading coefficients approximately (3, 0, -1, 1, 0, 1, -1, 0, 1, …) at
  50-digit precision (signed, using the centred variant). The expansion does not
  terminate, as pi notin Q(5).

  Phase transition exponent: muT(pi) = 0 (trivially, since pi is exactly representable as
  T(1,0,0,1,0) at complexity 4). This illustrates that elements of T are maximally
  Trinity-dominant.

  10.5 The AlphaPhi Disambiguation (flos38.tex)
  The PhD programme introduces an object called alpha$\varphi$ in Chapter 38, which appeared
  in early drafts with a notation inconsistency. This section provides the canonical
  disambiguation required for use throughout Paper 3 and in all future GOLDEN BRIDGE (Trinity-Pellis Bridge)
  papers.

  Two distinct objects:

  Object 1 --- Algebraic AlphaPhi:
  alpha$\varphi$( alg) = ($\varphi$-3)/(2) = ((5-2))/(2) approx 0.11803. 2
  This is a purely algebraic expression in Q(5). Its Pellis expansion terminates in exactly 4
  steps: alpha$\varphi$( alg) = $\varphi$-3/2, with Trinity label T(1, 0, -3, 0, 0)/2 (adjusting the rational
  prefactor). MDL score approx 7 bits. No logarithms appear. This object is
  Trinity-dominant with muT(alpha$\varphi$( alg)) = 0 (it lies in Q(5), the field of fractions of
  Z[$\varphi$]).

  Object 2 --- Logarithmic AlphaPhi:
  alpha$\varphi$(log) = (ln($\varphi$2))/(pi) = (2ln$\varphi$)/(pi) approx (0.96242)/(3.14159) approx 0.30635. 3

  This is a transcendental number (involving both ln$\varphi$ and pi). Its Trinity representation
  at low complexity: no monomial in T represents it exactly; the nearest at L = 10
  involves terms of the form $\varphi$p pim eq with large residuals. The Pellis expansion does not
  terminate (since alpha$\varphi$(log) notin Q(5)). Phase transition estimate: near the critical
  line, muT approx 1.0.

  These two objects are not equal: alpha$\varphi$( alg) approx 0.118 $\neq$ 0.306 approx alpha$\varphi$(log).

  The notation inconsistency in the early drafts of Chapter 38 apparently used the same
  symbol alpha$\varphi$ for both objects, leading to confusion in derived computations. Paper 3
  adopts the disambiguated notation above as the canonical definition for all subsequent
  work in this trilogy.

  Canonical roles in the Paper 3 formalism:

  10.6 The Fine-Structure Constant
  The fine-structure constant alpha QED approx 1/137.035999177 (NIST CODATA 2022)
  governs the strength of the electromagnetic interaction. It is a running coupling
  constant --- its value depends on the energy scale of the interaction --- and is
  determined by QED and the renormalisation group, not by any simple arithmetic
  formula.

  Mathematical context only. The framework asks: what is the most compact Trinity
  monomial approximation to alpha QED-1 approx 137.036?

  - T(137, 0, 0, 0, 0) = 137: MDL = 8 bits; relative error approx 2.6 $\times$ 10-4.
  - At L = 14: refined combinations such as T(137, 0, 0, 0, 0) + Pellis terms approximate
  alpha QED-1 to 3--4 significant figures.

  Pellis expansion: Leading coefficients approximately (137, 0, 0, 1, 0, -1, …) at 50-digit
  precision.

  Phase transition estimate: muT(alpha QED-1) approx 1.0 ± 0.1, near the critical line.
  The integer 137 provides a low-cost Trinity anchor; the residual approx 0.036 requires
  Pellis structure.

  Critical disclaimer. The numerical proximity of alpha QED-1 to 137 is an observed
  experimental fact, not a theoretical consequence of the Trinity framework. Neither the
  integer 137 nor any Trinity monomial is a theoretical explanation of the fine-structure
  constant's value. The MDL score does not distinguish genuine structural relationships
  from numerical coincidences; it is a necessary but not sufficient condition for the
  former.

  Disambiguation check. alpha$\varphi$( alg) approx 0.118 $\neq$ alpha QED approx 7.3 $\times$ 10-3;
  alpha$\varphi$(log) approx 0.306 $\neq$ alpha QED. No coincidence between the AlphaPhi objects
  and the QED coupling is asserted.

  10.7 The Proton-to-Electron Mass Ratio
  mupe = mp / me approx 1836.15267343 (NIST CODATA 2022). This ratio is determined
  by QCD, electroweak physics, and fundamental hadron structure --- not by any
  arithmetic formula.

  - T(1836, 0, 0, 0, 0) = 1836: relative error approx 8.3 $\times$ 10-5; MDL = 11 bits.
  - Pellis expansion of fractional part 0.15267: coefficients approximately (0, 0, 1, -1, 0, 1,
  …); slow convergence.

  Phase transition estimate: muT(mupe) approx 0.7 (Trinity-dominant). The integer 1836
  provides a good low-cost anchor; the Pellis expansion captures the remaining precision.

  Commentary. The Trinity-dominance of mupe at low complexity simply reflects the fact
  that 1836 is a good integer approximation to mupe. This says nothing about the
  physical origin of the mass ratio.

  10.8 Algebraic Constants: 2 and log 2
  2 approx 1.41421: The illustrative Trinity approximant T(1,0,1,-1,1) = $\varphi$ e/pi approx
  1.40001 has relative error approx 1.0\%. Since 2 notin Q(5), its Pellis expansion is
  aperiodic. Phase transition estimate: muT(2) approx 1.3 (slightly Pellis-dominant), as
  expected for an algebraic number not in Q(5).

  log 2 approx 0.69315: The nearest Trinity approximant at L = 10 is T(1,-1,1,0,0) = $\varphi$/3
  approx 0.539 (relative error 22\%). Algebraic numbers not involving $\varphi$ do not compress
  well under Trinity representation --- an expected feature that confirms the framework's
  $\varphi$-centric design.

  10.9 Fibonacci/Quasicrystal Connection
  A key motivation for the Trinity framework is the discovery by Shechtman et al. (1984)
  of quasicrystalline materials with icosahedral symmetry, in which the ratio of
  interatomic distances is governed by $\varphi$ and the structure is indexed by Fibonacci
  sequences. As the Nobel Prize scientific background on quasicrystals notes, in Penrose
  tilings and icosahedral quasicrystals, "the ratio of various distances between atoms is
  always related to tau [the golden ratio]" and "both the Fibonacci sequence and the
  golden ratio are important to scientists when they want to use a diffraction pattern to
  describe quasicrystals at the atomic level."

  In the Trinity framework, this manifests as follows. The Fibonacci numbers Fn and
  Lucas numbers Ln are the canonical elements of Z[$\varphi$]; their Pellis expansions terminate
  exactly. A quasicrystalline material indexed by Fibonacci-number spacings
  corresponds, in the Trinity framework, to a physical structure whose geometric ratios
  are exactly representable in Z[$\varphi$] --- maximally Trinity-dominant (Proposition 8.2). The
  aperiodic but ordered structure of quasicrystals is thus a physical realisation of
  Trinity-dominance.

  This analogy is mathematical and conceptual. The Paper 3 framework does not model
  the physics of quasicrystals, and no claim is made about physical structure from the
  symbolic framework alone.

\subsection{References}

- Vasilev, D. \textit{Flos Aureus Monograph: Trinity Constants}. DOI \href{https://zenodo.org/doi/10.5281/zenodo.19227877}{10.5281/zenodo.19227877}.
- CODATA 2018 recommended values: NIST \href{https://physics.nist.gov/cuu/Constants/}{https://physics.nist.gov/cuu/Constants/}.
- Heyrovska, R. Golden-angle FSC publications --- Heyrovský Institute of Physical Chemistry.
